En una ciutat, el 70\,\% de les persones llegeix el diari digital ($D$), el 40\,\% escolta la ràdio ($R$) i
el 20\,\% no fa cap de les dues coses. Es tria una persona a l'atzar.

\begin{apartats}

\apartat[1,25]{0,75}
Quina és la probabilitat que faci totes dues coses?

\begin{solucio}
$P(D\cup R)=1-0{,}2=0{,}8$, perquè «no fer-ne cap» és el contrari de la unió.\\
$P(D\cap R)=P(D)+P(R)-P(D\cup R)=0{,}7+0{,}4-0{,}8=0{,}3$.
\end{solucio}

\begin{tria}{operacions-tasca}
\itemtria{original}{1}{1,25}
Quina és la probabilitat que només escolti la ràdio? I que faci només una de les dues coses?

\begin{solucio}
\textbf{Només ràdio}: $P\left(R\cap\overline{D}\right)=P(R)-P(D\cap R)=0{,}4-0{,}3=0{,}1$.\\
\textbf{Només una}: $P(D\cup R)-P(D\cap R)=0{,}8-0{,}3=0{,}5$.
\end{solucio}

\itemtria{dades-impossibles}{1}{1,25}
Una altra enquesta diu que el 70\,\% de les persones llegeix el diari digital, el 40\,\% escolta la ràdio
i només el 5\,\% fa totes dues coses. Són possibles aquestes dades? Justifica-ho.

\begin{solucio}
\textbf{No}: donarien $P(D\cup R)=0{,}7+0{,}4-0{,}05=1{,}05>1$.\\
Com que $P(D\cup R)\le1$, ha de ser $P(D\cap R)\ge0{,}7+0{,}4-1=0{,}1$: com a mínim, el 10\,\% hauria de
fer totes dues coses.
\end{solucio}
\end{tria}

\begin{nomesllarg}
\apartat{0,75}
Si una persona escolta la ràdio, quina és la probabilitat que també llegeixi el diari digital?

\begin{solucio}
$P(D\mid R)=\dfrac{P(D\cap R)}{P(R)}=\dfrac{0{,}3}{0{,}4}=0{,}75$.
\end{solucio}
\end{nomesllarg}

\end{apartats}
